\documentclass[a4paper,twoside]{article}

\usepackage{morewrites}

\usepackage[svgnames,dvipsnames,x11names]{xcolor}
\usepackage{graphicx}
\usepackage[english]{babel}
\usepackage[colorlinks=true, linkcolor=NavyBlue, urlcolor=NavyBlue, citecolor=NavyBlue]{hyperref}
\usepackage[a4paper]{geometry}
\usepackage{ts-fonts}
\usepackage{csquotes}
\usepackage{amsmath}
\usepackage{float}
\usepackage{seqsplit}
\usepackage{ts-code}

\usepackage{lastpage}
\usepackage{refcount}
\makeatletter
\def\ps@plain{
  \let\@oddhead\@empty
  \let\@evenhead\@empty
  \def\@oddfoot{
    \hfil
    \ifnum\getpagerefnumber{LastPage}>1\relax
      \thepage
    \fi
    \hfil}
  \let\@evenfoot\@oddfoot
}
\makeatother

\title{Managing Slots}
\author{}\date{}

\begin{document}

\maketitle

\thispagestyle{plain}
Slots are placeholders in the \LaTeX{} template where specific document sections can be injected. Use the \texttt{-\allowbreak{}-\allowbreak{}slot} (or \texttt{-\allowbreak{}s}) option to map input documents to these slots. For example, to inject \texttt{abstract.md} into the \texttt{abstract} slot and \texttt{dedication.md} into the \texttt{dedication} slot of a book template, run:

\begin{code}{bash}{}{}
\begin{Verbatim}[commandchars=\\\{\},breaklines, breakanywhere, commandchars=\\\{\}]
texsmith\PY{+w}{ }abstract.md\PY{+w}{ }dedication.md\PY{+w}{ }chapter*.md\PY{+w}{ }\PY{l+s+se}{\PYZbs{}}
\PY{+w}{  }\PYZhy{}\PYZhy{}template\PY{+w}{ }book\PY{+w}{ }\PY{l+s+se}{\PYZbs{}}
\PY{+w}{  }\PYZhy{}\PYZhy{}slot\PY{+w}{ }abstract:abstract.md\PY{+w}{ }\PY{l+s+se}{\PYZbs{}}
\PY{+w}{  }\PYZhy{}\PYZhy{}slot\PY{+w}{ }dedication:dedication.md
\end{Verbatim}
\end{code}
To see the available slots for a given template, use the \texttt{-\allowbreak{}-\allowbreak{}template-\allowbreak{}info} flag:

\begin{code}{text}{}{baselinestretch=0.5, }
\begin{Verbatim}[commandchars=\\\{\},breaklines, breakanywhere, commandchars=\\\{\}]
\PYZdl{} uv run texsmith \PYZhy{}tbook \PYZhy{}\PYZhy{}template\PYZhy{}info
...
                                     Slots
┌────────────┬─────────┬───────────┬─────────┬────────┬────────────┬───────────┐
│            │         │      Base │         │        │  Effective │   Strip   │
│ Name       │ Default │     Level │   Depth │ Offset │      Level │  Heading  │
├────────────┼─────────┼───────────┼─────────┼────────┼────────────┼───────────┤
│ appendix   │         │         \PYZhy{} │ chapter │      0 │          0 │    no     │
│ backmatter │         │         \PYZhy{} │ chapter │      0 │          0 │    no     │
│ colophon   │         │         \PYZhy{} │ chapter │      0 │          0 │    yes    │
│ dedication │         │         \PYZhy{} │ chapter │      0 │          0 │    yes    │
│ frontmatt… │         │         \PYZhy{} │ chapter │      0 │          0 │    no     │
│ mainmatter │    *    │         \PYZhy{} │ chapter │      0 │          0 │    no     │
│ preface    │         │         0 │       \PYZhy{} │      0 │          0 │    no     │
└────────────┴─────────┴───────────┴─────────┴────────┴────────────┴───────────┘
\end{Verbatim}
\end{code}
You can also extract sections from a single document using the \texttt{slot:Section Name} syntax. For example, to inject the \enquote{Abstract} section from \texttt{main.md} into the \texttt{abstract} slot:

\begin{code}{bash}{}{}
\begin{Verbatim}[commandchars=\\\{\},breaklines, breakanywhere, commandchars=\\\{\}]
texsmith\PY{+w}{ }main.md\PY{+w}{ }\PY{l+s+se}{\PYZbs{}}
\PY{+w}{  }\PYZhy{}\PYZhy{}template\PY{+w}{ }article\PY{+w}{ }\PY{l+s+se}{\PYZbs{}}
\PY{+w}{  }\PYZhy{}\PYZhy{}slot\PY{+w}{ }abstract:slot:Abstract
\end{Verbatim}
\end{code}

\end{document}
